\section{List of workflow tasks}

This table lists all the tasks and associated recipes in the \instname\ workflow.
Only some tasks are involved in the scientific reduction. These are marked as
``yes'' (executed by default) or ``optional'' (executed only when required by
the workflow). Tasks marked as ``no'' are not part of the default science
reduction; they are intended primarily for instrument monitoring and can be
executed only when explicitly selected as the reduction target. When a task is
selected as the target, all prerequisite tasks required to generate its input
calibrations are executed automatically.

\begin{table}[htbp]
\centering
\scriptsize
\begin{tabular}{|l|l|c|p{7cm}|}
\hline
\textbf{TASK} & \textbf{RECIPE} & \textbf{Used in science} & \textbf{Notes} \\
              &                 & \textbf{reduction} & \\
\hline
ifs\_coronagraph\_center &
sph\_ifs\_science\_dr &
optional &
Determines the stellar position behind the coronagraph. \\
\hline
ifs\_dark &
sph\_ifs\_master\_dark &
yes &
Creates the master dark (\texttt{IFS\_MASTER\_DARK}). \\
\hline
ifs\_det\_flat\_broad\_band &
sph\_ifs\_master\_detector\_flat &
yes &
Generates the broadband detector flat used for small-scale pixel-to-pixel response correction. \\
\hline
ifs\_det\_flat\_narrow\_band1 &
sph\_ifs\_master\_detector\_flat &
yes &
Produces the first narrow-band detector flat component. \\
\hline
ifs\_det\_flat\_narrow\_band2 &
sph\_ifs\_master\_detector\_flat &
yes &
Produces the second narrow-band detector flat component. \\
\hline
ifs\_det\_flat\_narrow\_band3 &
sph\_ifs\_master\_detector\_flat &
yes &
Produces the third narrow-band detector flat component. \\
\hline
ifs\_det\_flat\_narrow\_band4 &
sph\_ifs\_master\_detector\_flat &
yes &
Produces the fourth narrow-band detector flat component (associated only for the \texttt{PRI\_YJH} disperser). \\
\hline
ifs\_distortion\_map &
sph\_ifs\_distortion\_map &
no &
Measures the geometric distortion of the IFS lenslet grid. Produces \texttt{IFS\_DISTORTION\_MAP} and, when needed, \texttt{IFS\_POINT\_PATTERN}. \\
\hline
ifs\_ifu\_flat &
sph\_ifs\_instrument\_flat &
yes &
Produces the IFU (lenslet) flat-field calibration used by subsequent reductions. \\
\hline
ifs\_instrument\_background &
sph\_ifs\_cal\_background &
yes &
Produces the instrumental background calibration. \\
\hline
ifs\_science &
sph\_ifs\_science\_dr &
yes &
Main science reduction task applying background subtraction, flat-fielding, wavelength calibration and optional ADI or spectral deconvolution. \\
\hline
ifs\_science\_flux &
sph\_ifs\_science\_dr &
optional &
Processes FLUX observations obtained with the target offset from the coronagraph to measure the stellar flux. \\
\hline
ifs\_sky\_background &
sph\_ifs\_cal\_background &
yes &
Produces the sky-background calibration. \\
\hline
ifs\_spectra\_positions &
sph\_ifs\_spectra\_positions &
yes &
Associates detector pixels with lenslets and generates the initial Pixel Description Table (\texttt{IFS\_SPECPOS}). \\
\hline
ifs\_standard\_astrometry &
sph\_ifs\_science\_dr &
optional &
Processes astrometric standard observations. \\
\hline
ifs\_standard\_flux &
sph\_ifs\_science\_dr &
optional &
Processes spectrophotometric standard observations. \\
\hline
ifs\_static\_bad\_pixel\_map &
sph\_ifs\_master\_dark &
no &
Produces the static bad-pixel map (\texttt{IFS\_STATIC\_BADPIXELMAP}). \\
\hline
ifs\_wavelength\_calibration &
sph\_ifs\_wave\_calib &
yes &
Refines the wavelength calibration and produces \texttt{IFS\_WAVECALIB}. \\
\hline
\end{tabular}
\caption{Tasks and associated recipes of the \instname\ workflow.}
\label{tab:ifs_workflow_tasks}
\end{table}


\paragraph{Notes}

\begin{itemize}

\item Background calibrations are associated in the following priority order:
\texttt{SKY\_BG}
$\rightarrow$
\texttt{INS\_BG}
$\rightarrow$
\texttt{MASTER\_DARK}.

\item The science reduction is performed on a file-by-file basis. The pipeline does not combine multiple science exposures; any stacking or post-processing of multiple reduced frames must be performed afterwards.

\item The \texttt{force\_distortion\_correction} workflow parameter is set to
\texttt{FALSE} by default for science and IDP reductions (see Section~\ref{sec:distortion_map}). Although a distortion calibration may be associated by EDPS, it is not propagated to the science recipe. For QC0 processing, this parameter is set to \texttt{TRUE}.

\end{itemize}